\chapter{Serii Taylor}

\section{Exerciții suplimentare}

1. Găsiți aproximarea liniară și pătratică a funcțiilor:
\begin{enumerate}[(a)]
\item $ f(x) = \sqrt[3]{x} $;
\item $ f(x) = \sin(\cos x) $;
\item $ f(x) = e^{\sin x} $;
\item $ f(x) = \arcsin x $.
\end{enumerate}

\vspace{1cm}

2. Folosind seria Taylor, aproximați cu o eroare mai mică decît $ 10^{-3} $
numerele:
\begin{enumerate}[(a)]
\item $ \sqrt[3]{65} $;
\item $ \sin 32 $;
\item $ \arctan \dfrac{1}{2} $;
\item $ e^{-0,2} $;
\item $ \ln 1,1 $;
\item $ \ln 4 $;
\item $ \ln 5 $.
\end{enumerate}

\emph{Indicație:} Atenție la domeniile de convergență!

\vspace{1cm}

3. Să se calculeze cu o eroare mai mică decît $ 10^{-3} $ integralele:
\begin{enumerate}[(a)]
\item $ \ds\int_0^{\frac{1}{2}} \dfrac{\sin x}{x} dx $;
\item $ \ds\int_0^{\frac{1}{2}} \dfrac{\ln(1 + x)}{x} dx $;
\item $ \ds\int_0^{\frac{1}{3}} \dfrac{\arctan x}{x} dx $;
\item $ \ds\int_0^1 e^{-x^2} dx $.
\end{enumerate}

\vspace{1cm}

4.  Găsiți mulțimea de convergență și suma seriei:
\[
  sum_{n \geq 0} (-1)^n \dfrac{x^{2n + 1}}{2n + 1}.
\]

\emph{Indicație:} $ R = 1 $ (raport), iar suma se poate afla derivînd
termen cu termen. Rezultă (prin derivare) seria geometrică de rază $ -x^2 $,
cu suma $ \dfrac{1}{x^2 + 1} $, valabilă pentru $ x \in (-1, 1) $.

Rezultă $ f(x) = \arctan x + c $.

\vspace{1cm}

5. Arătați că seriile numerice de mai jos sînt convergente și calculați
sumele lor, folosind serii de puteri:
\begin{enumerate}[(a)]
\item $ \ds\sum_{n \geq 0} \dfrac{(-1)^n}{3n + 1} $;
\item $ \ds\sum_{n \geq 0} \dfrac{(n+1)^2}{n!} $;
\item $ \ds\sum_{n \geq 1} \dfrac{n^2(3^n - 2^n)}{6^n} $.
\end{enumerate}

\emph{Indicații:}
\begin{enumerate}[(a)]
\item Seria satisface criteriul lui Leibniz, deci este convergentă.

  Pentru a găsi suma, pornim cu seria de puteri $ \sum (-1)^n \dfrac{x^{3n+1}}{3n+1} $.

  Intervalul de convergență este $ (-1, 1) $, iar pentru $ x = 1 $, avem
  seria dată.

  Fie $ f $ suma acestei serii de puteri în intervalul $ (-1, 1) $. Derivăm
  termen cu termen și obținem:
  \[
    f'(x) = \sum (-1)^n x^{3n} = \frac{1}{1 + x^3},
  \]
  pentru $ |x| < 1 $, ca suma unei serii geometrice alternate.

  Rezultă:
  \[
    f(x) = \int \frac{dx}{1 + x^3} = \frac{1}{6} \ln \frac{(x + 1)^2}{x^2 - x + 1} %
    + \frac{1}{\sqrt{3}} \arctan{2x - 1}{\sqrt{3}} + c,
  \]
  pentru $ |x| < 1 $. Calculînd $ f(0) $, găsim $ c = \dfrac{\pi}{6 \sqrt{3}} $.
\item Se folosește seria pentru $ e^x $, din care obținem seria pentru
  $ (x + x^2)e^x $, pe care o derivăm termen cu termen.

  Pentru $ x = 1 $, se obține seria cerută, cu suma $ 5e $.
\item Descompunem seria în două, apoi folosim seria de puteri $ \sum n^2 x^n $,
  pe care o derivăm termen cu termen, pentru a obține seria pentru $ nx^{n-1} $,
  apoi seria pentru $ nx^n $.
\end{enumerate}




%%% Local Variables:
%%% mode: latex
%%% TeX-master: "../m1-18-19"
%%% End:
